% Options for packages loaded elsewhere
\PassOptionsToPackage{unicode}{hyperref}
\PassOptionsToPackage{hyphens}{url}
%
\documentclass[
  11pt,
]{article}
\usepackage{amsmath,amssymb}
\usepackage{iftex}
\ifPDFTeX
  \usepackage[T1]{fontenc}
  \usepackage[utf8]{inputenc}
  \usepackage{textcomp} % provide euro and other symbols
\else % if luatex or xetex
  \usepackage{unicode-math} % this also loads fontspec
  \defaultfontfeatures{Scale=MatchLowercase}
  \defaultfontfeatures[\rmfamily]{Ligatures=TeX,Scale=1}
\fi
\usepackage{lmodern}
\ifPDFTeX\else
  % xetex/luatex font selection
\fi
% Use upquote if available, for straight quotes in verbatim environments
\IfFileExists{upquote.sty}{\usepackage{upquote}}{}
\IfFileExists{microtype.sty}{% use microtype if available
  \usepackage[]{microtype}
  \UseMicrotypeSet[protrusion]{basicmath} % disable protrusion for tt fonts
}{}
\makeatletter
\@ifundefined{KOMAClassName}{% if non-KOMA class
  \IfFileExists{parskip.sty}{%
    \usepackage{parskip}
  }{% else
    \setlength{\parindent}{0pt}
    \setlength{\parskip}{6pt plus 2pt minus 1pt}}
}{% if KOMA class
  \KOMAoptions{parskip=half}}
\makeatother
\usepackage{xcolor}
\usepackage[margin=2.2cm]{geometry}
\setlength{\emergencystretch}{3em} % prevent overfull lines
\providecommand{\tightlist}{%
  \setlength{\itemsep}{0pt}\setlength{\parskip}{0pt}}
\setcounter{secnumdepth}{-\maxdimen} % remove section numbering
\ifLuaTeX
  \usepackage{selnolig}  % disable illegal ligatures
\fi
\IfFileExists{bookmark.sty}{\usepackage{bookmark}}{\usepackage{hyperref}}
\IfFileExists{xurl.sty}{\usepackage{xurl}}{} % add URL line breaks if available
\urlstyle{same}
\hypersetup{
  pdftitle={Worksheet 5E: Laws of Sines and Cosines},
  hidelinks,
  pdfcreator={LaTeX via pandoc}}

\title{Worksheet 5E: Laws of Sines and Cosines}
\author{}
\date{}

\begin{document}
\maketitle

\begin{quote}
\textbf{What This Worksheet Is About}

Right triangle trigonometry only works for\ldots{} right triangles. But
most triangles in the real world aren't right triangles!

The Law of Sines and Law of Cosines extend trigonometry to ALL
triangles. These tools are essential for surveying, navigation,
astronomy, and engineering.
\end{quote}

\begin{center}\rule{0.5\linewidth}{0.5pt}\end{center}

\hypertarget{part-a-the-law-of-sines-22-problems}{%
\subsection{Part A: The Law of Sines (22
problems)}\label{part-a-the-law-of-sines-22-problems}}

\hypertarget{the-formula}{%
\subsubsection{The Formula}\label{the-formula}}

For ANY triangle with sides \(a\), \(b\), \(c\) opposite angles \(A\),
\(B\), \(C\):

\[\boxed{\frac{a}{\sin A} = \frac{b}{\sin B} = \frac{c}{\sin C}}\]

Equivalently: \(\frac{\sin A}{a} = \frac{\sin B}{b} = \frac{\sin C}{c}\)

\textbf{When to use:} When you know an angle and its opposite side, plus
one other piece.

\begin{center}\rule{0.5\linewidth}{0.5pt}\end{center}

\textbf{Find the missing parts (round to 1 decimal place):}

\textbf{1.} Triangle with \(A = 40°\), \(B = 60°\), \(a = 10\) cm. Find
\(b\) and \(c\).

\emph{Hint: First find \(C = 180° - A - B\)}

\textbf{2.} Triangle with \(A = 35°\), \(C = 75°\), \(c = 15\) m. Find
\(a\) and \(b\).

\textbf{3.} Triangle with \(B = 52°\), \(C = 48°\), \(b = 8\) cm. Find
\(a\) and \(c\).

\textbf{4.} Triangle with \(A = 100°\), \(B = 35°\), \(a = 20\) m. Find
\(b\) and \(c\).

\begin{center}\rule{0.5\linewidth}{0.5pt}\end{center}

\hypertarget{the-ambiguous-case-ssa}{%
\subsubsection{The Ambiguous Case (SSA)}\label{the-ambiguous-case-ssa}}

When given two sides and an angle NOT between them (SSA), there may be:
- \textbf{No solution} (triangle impossible) - \textbf{One solution}
(unique triangle) - \textbf{Two solutions} (two different triangles
work!)

This happens because \(\sin\theta = \sin(180° - \theta)\).

\begin{center}\rule{0.5\linewidth}{0.5pt}\end{center}

\textbf{5.} Triangle with \(A = 30°\), \(a = 5\), \(b = 8\).

Find angle \(B\): \(\frac{\sin B}{8} = \frac{\sin 30°}{5}\)

\(\sin B = \frac{8 \cdot 0.5}{5} = 0.8\)

\(B = \sin^{-1}(0.8) = ?\) or \(B = 180° - ?\)

Are both solutions valid? (Check if \(A + B < 180°\))

\textbf{6.} Triangle with \(A = 40°\), \(a = 12\), \(b = 18\). How many
solutions?

\textbf{7.} Triangle with \(A = 60°\), \(a = 10\), \(b = 8\). How many
solutions?

\textbf{8.} Triangle with \(A = 45°\), \(a = 6\), \(b = 10\). How many
solutions?

\begin{center}\rule{0.5\linewidth}{0.5pt}\end{center}

** Astronomy: Stellar Triangulation**

\textbf{9.} The parallax method for measuring star distances creates a
triangle with: - Earth at two positions 6 months apart (baseline = 2 AU
≈ 300 million km) - The star at the third vertex - The parallax angle
\(p\) at the star

For Sirius, if the angle at one Earth position is 89.999623° and at the
other is 89.999631°: - (a) What is the angle at Sirius (the parallax
angle)? - (b) Using Law of Sines with baseline 300 million km, estimate
the distance to Sirius.

\emph{Note: Actual parallax measurements are much more precise than this
simplified example!}

\textbf{10.} Alpha Centauri has a parallax of about 0.76 arcseconds
(convert to degrees: 1° = 3600 arcseconds).

Set up the Law of Sines equation to find its distance from the Sun.

\begin{center}\rule{0.5\linewidth}{0.5pt}\end{center}

** Microscopy: Electron Beam Triangulation**

\textbf{11.} In some electron microscope techniques, the beam hits a
sample at an angle and the scattered electrons are detected.

If the beam hits at 15° from vertical, travels to the sample, and
electrons scatter to a detector at 25° from vertical (on the opposite
side), with the sample-to-detector distance being 50 mm:

Set up a triangle and use Law of Sines to find the beam path length from
source to sample.

\begin{center}\rule{0.5\linewidth}{0.5pt}\end{center}

** Surveying: Triangulation**

\textbf{12.} Surveyors measure the distance between two points A and B
as 500 m. From A, the angle to a distant point C is 72°. From B, the
angle to C is 63°.

\begin{itemize}
\item
  \begin{enumerate}
  \def\labelenumi{(\alph{enumi})}
  \tightlist
  \item
    What is angle C in the triangle ABC?
  \end{enumerate}
\item
  \begin{enumerate}
  \def\labelenumi{(\alph{enumi})}
  \setcounter{enumi}{1}
  \tightlist
  \item
    Find the distances AC and BC.
  \end{enumerate}
\end{itemize}

\textbf{13.} From two observation points 2 km apart, a mountain peak is
sighted. The angles of elevation aren't needed --- just the horizontal
angles: from point A, the bearing to the peak is N 35° E, and from point
B (due east of A), the bearing is N 52° W.

\begin{itemize}
\item
  \begin{enumerate}
  \def\labelenumi{(\alph{enumi})}
  \tightlist
  \item
    Sketch the triangle and find all three angles.
  \end{enumerate}
\item
  \begin{enumerate}
  \def\labelenumi{(\alph{enumi})}
  \setcounter{enumi}{1}
  \tightlist
  \item
    How far is the peak from point A?
  \end{enumerate}
\end{itemize}

\begin{center}\rule{0.5\linewidth}{0.5pt}\end{center}

\hypertarget{part-b-the-law-of-cosines-22-problems}{%
\subsection{Part B: The Law of Cosines (22
problems)}\label{part-b-the-law-of-cosines-22-problems}}

\hypertarget{the-formula-1}{%
\subsubsection{The Formula}\label{the-formula-1}}

For ANY triangle with sides \(a\), \(b\), \(c\) opposite angles \(A\),
\(B\), \(C\):

\[\boxed{c^2 = a^2 + b^2 - 2ab\cos C}\]

Or equivalently for any side: - \(a^2 = b^2 + c^2 - 2bc\cos A\) -
\(b^2 = a^2 + c^2 - 2ac\cos B\)

\textbf{When to use:} - SAS: Two sides and the included angle - SSS: All
three sides (to find angles)

\textbf{Note:} When \(C = 90°\), this becomes Pythagoras! (since
\(\cos 90° = 0\))

\begin{center}\rule{0.5\linewidth}{0.5pt}\end{center}

\textbf{Find the missing side:}

\textbf{14.} \(a = 7\), \(b = 10\), \(C = 50°\). Find \(c\).

\textbf{15.} \(b = 12\), \(c = 8\), \(A = 110°\). Find \(a\).

\textbf{16.} \(a = 5\), \(c = 9\), \(B = 72°\). Find \(b\).

\textbf{17.} \(a = 15\), \(b = 20\), \(C = 35°\). Find \(c\).

\begin{center}\rule{0.5\linewidth}{0.5pt}\end{center}

\textbf{Find the missing angle:}

\textbf{18.} \(a = 5\), \(b = 7\), \(c = 9\). Find angle \(C\).

\textbf{19.} \(a = 8\), \(b = 8\), \(c = 10\). Find angle \(C\).

\textbf{20.} \(a = 6\), \(b = 10\), \(c = 12\). Find angle \(A\).

\textbf{21.} \(a = 13\), \(b = 14\), \(c = 15\). Find all three angles.

\begin{center}\rule{0.5\linewidth}{0.5pt}\end{center}

** Astronomy: Planetary Distances**

\textbf{22.} At a certain moment: - Earth is 1 AU from the Sun - Mars is
1.5 AU from the Sun\\
- The angle at the Sun between Earth and Mars is 60°

Find the distance from Earth to Mars at this moment.

\textbf{23.} Venus is 0.72 AU from the Sun. When Venus is at ``greatest
elongation'' (maximum angular distance from the Sun as seen from Earth),
the Earth-Venus-Sun angle is 90°.

Using the Law of Cosines with Earth at 1 AU and the Sun-Earth-Venus
angle at greatest elongation: - (a) If the Sun-Earth-Venus angle is 46°,
find the Earth-Venus distance.

\begin{center}\rule{0.5\linewidth}{0.5pt}\end{center}

** Crystallography: Interatomic Distances**

\textbf{24.} In a crystal, three atoms form a triangle. Atom A is bonded
to atom B (distance 0.154 nm) and to atom C (distance 0.143 nm). The
bond angle at A is 109.5° (tetrahedral angle).

Find the distance from B to C.

\textbf{25.} A water molecule has two O-H bonds of length 0.096 nm each,
with a bond angle of 104.5°.

Find the distance between the two hydrogen atoms.

\begin{center}\rule{0.5\linewidth}{0.5pt}\end{center}

** Microscopy: Stage Geometry**

\textbf{26.} A microscope sample stage can translate and rotate. A
feature on the sample is at position A. After translating 5 mm and
rotating 30°, it appears at position B relative to the objective.

If the rotation is about a point 10 mm from the original feature
position: - Model this as a triangle with the rotation center, original
position, and final position. - Find how far the feature has actually
moved (distance from A to B).

\begin{center}\rule{0.5\linewidth}{0.5pt}\end{center}

** Space Navigation**

\textbf{27.} A spacecraft needs to travel from point A to point B, but
must avoid a debris field. It goes from A to waypoint C (distance 500
km), then from C to B (distance 700 km). The angle at C is 140°.

What is the direct distance from A to B that the spacecraft is avoiding?

\textbf{28.} A satellite in orbit has position vectors at two times: -
Time 1: 7000 km from Earth center, at angle 0° - Time 2: 7500 km from
Earth center, at angle 25° (angular separation)

Find the distance traveled along the orbit (straight-line distance, not
arc length).

\begin{center}\rule{0.5\linewidth}{0.5pt}\end{center}

\hypertarget{part-c-area-of-triangles-14-problems}{%
\subsection{Part C: Area of Triangles (14
problems)}\label{part-c-area-of-triangles-14-problems}}

\hypertarget{multiple-formulas}{%
\subsubsection{Multiple Formulas}\label{multiple-formulas}}

\textbf{Base and Height:} \[A = \frac{1}{2}bh\]

\textbf{Two Sides and Included Angle:}
\[\boxed{A = \frac{1}{2}ab\sin C}\]

\textbf{Heron's Formula (three sides):} \[A = \sqrt{s(s-a)(s-b)(s-c)}\]
where \(s = \frac{a+b+c}{2}\) (semi-perimeter)

\begin{center}\rule{0.5\linewidth}{0.5pt}\end{center}

\textbf{Find the area:}

\textbf{29.} Triangle with \(a = 8\), \(b = 12\), \(C = 40°\).

\textbf{30.} Triangle with \(b = 15\), \(c = 20\), \(A = 70°\).

\textbf{31.} Triangle with sides 5, 12, 13.

\textbf{32.} Triangle with sides 7, 8, 9.

\textbf{33.} Triangle with \(A = 50°\), \(B = 60°\), \(c = 10\).

\emph{Hint: Use Law of Sines to find other sides first.}

\begin{center}\rule{0.5\linewidth}{0.5pt}\end{center}

** Science Connection: Cross Products**

\textbf{34.} In physics, the cross product of two vectors \(\vec{u}\)
and \(\vec{v}\) has magnitude:
\[|\vec{u} \times \vec{v}| = |\vec{u}||\vec{v}|\sin\theta\]

This equals twice the area of the triangle formed by the vectors!

If \(|\vec{u}| = 5\), \(|\vec{v}| = 8\), and \(\theta = 30°\), find: -
(a) The magnitude of the cross product - (b) The area of the triangle
formed by the vectors

\begin{center}\rule{0.5\linewidth}{0.5pt}\end{center}

** Astronomy: Orbital Sectors**

\textbf{35.} Kepler's Second Law says a planet sweeps out equal areas in
equal times.

If Earth moves through 1° of its orbit in about 1 day (simplification),
and Earth's distance from the Sun varies from 147 million km
(perihelion) to 152 million km (aphelion):

\begin{itemize}
\item
  \begin{enumerate}
  \def\labelenumi{(\alph{enumi})}
  \tightlist
  \item
    Calculate the area swept in one day at perihelion.
  \end{enumerate}
\item
  \begin{enumerate}
  \def\labelenumi{(\alph{enumi})}
  \setcounter{enumi}{1}
  \tightlist
  \item
    Calculate the area swept in one day at aphelion.
  \end{enumerate}
\item
  \begin{enumerate}
  \def\labelenumi{(\alph{enumi})}
  \setcounter{enumi}{2}
  \tightlist
  \item
    If the areas are equal (Kepler's law), what does this tell you about
    Earth's speed at perihelion vs.~aphelion?
  \end{enumerate}
\end{itemize}

\begin{center}\rule{0.5\linewidth}{0.5pt}\end{center}

** Lithography: Pattern Area**

\textbf{36.} A triangular pattern is being etched onto a chip. The
triangle has vertices determined by: - Side 1: 500 nm - Side 2: 400 nm\\
- Angle between them: 60°

Find the area of this triangular feature in nm².

\begin{center}\rule{0.5\linewidth}{0.5pt}\end{center}

\hypertarget{part-d-applications-and-problem-solving-16-problems}{%
\subsection{Part D: Applications and Problem Solving (16
problems)}\label{part-d-applications-and-problem-solving-16-problems}}

** Astronomy: Lunar Distance**

\textbf{37.} Ancient astronomers estimated the Moon's distance using a
lunar eclipse. When the Moon passes through Earth's shadow, the geometry
creates triangles.

If the shadow at the Moon's distance has apparent angular size 1.4° (as
seen from Earth's center), and this shadow is created by the Sun
(angular size 0.53° from Earth), and the Sun is 150 million km away:

This is a complex problem! Start by drawing the geometry and identifying
the triangles involved.

\begin{center}\rule{0.5\linewidth}{0.5pt}\end{center}

** X-ray Crystallography**

\textbf{38.} In analyzing a protein crystal, X-ray diffraction reveals
that three atoms form a triangle with: - Side AB: determined to be 0.28
nm - Angle at A: 108° - Angle at B: 34°

Find: - (a) Angle at C - (b) Side AC - (c) Side BC - (d) The area of the
triangle

\begin{center}\rule{0.5\linewidth}{0.5pt}\end{center}

** Architecture: Geodesic Dome**

\textbf{39.} A geodesic dome is made of triangular panels. One panel
has: - Side lengths: 2.1 m, 2.1 m, 1.8 m (isosceles triangle)

Find: - (a) All three angles - (b) The area of one panel - (c) If the
dome has 320 such panels, what is the total surface area?

\begin{center}\rule{0.5\linewidth}{0.5pt}\end{center}

** Electron Diffraction**

\textbf{40.} When electrons pass through a crystal, they diffract
according to the crystal structure. The diffraction pattern reveals
interatomic distances.

A triangle of atoms in graphene has: - All sides equal to 0.142 nm
(equilateral)

Find: - (a) All angles (you already know this!) - (b) The area of this
triangular cell - (c) The height of the triangle

\begin{center}\rule{0.5\linewidth}{0.5pt}\end{center}

** GPS Trilateration**

\textbf{41.} GPS uses trilateration (circles) rather than triangulation
(angles), but the math involves similar geometry.

Three satellites report distances to a receiver: - Satellite A: 22,000
km - Satellite B: 23,000 km - Satellite C: 21,500 km

The satellites' positions form a triangle in space. If the angle at the
receiver between satellites A and B is 35°:

Use the Law of Cosines to find the distance from satellite A to
satellite B.

\begin{center}\rule{0.5\linewidth}{0.5pt}\end{center}

** Atomic Force Microscopy**

\textbf{42.} An AFM tip is positioned using piezoelectric actuators that
move in x, y, and z. When scanning at an angle to avoid obstacles:

The tip moves 2 μm in one direction, rotates, then moves 3 μm in a
direction 75° from the first.

Find: - (a) The straight-line displacement from start to finish - (b)
The angle this displacement makes with the first direction

\begin{center}\rule{0.5\linewidth}{0.5pt}\end{center}

** Binary Star Systems**

\textbf{43.} A binary star system has two stars orbiting their common
center of mass. At one instant: - Star A is 5 AU from the center of mass
- Star B is 3 AU from the center of mass - The angle at the center of
mass between the stars is 180° (opposite sides)

What is the distance between the two stars? (Easy!)

At another instant, the angle is 120°. Now what is the distance?

\begin{center}\rule{0.5\linewidth}{0.5pt}\end{center}

** Scanning Electron Microscope Geometry**

\textbf{44.} In an SEM, the electron beam, sample surface, and detector
form a triangle. - Beam to sample: 10 mm (working distance) - Sample to
detector: 30 mm - Beam angle from vertical: 0° (straight down) -
Detector angle from sample surface: 35°

Find the beam-to-detector distance.

\emph{Hint: You need to figure out the angle at the sample first.}

\begin{center}\rule{0.5\linewidth}{0.5pt}\end{center}

** Orbital Rendezvous**

\textbf{45.} Two spacecraft are in different orbits. At a certain time:
- Spacecraft A is 400 km above Earth (Earth radius ≈ 6371 km) -
Spacecraft B is 800 km above Earth - The angle at Earth's center between
them is 10°

Find the distance between the spacecraft.

\begin{center}\rule{0.5\linewidth}{0.5pt}\end{center}

\textbf{Final Challenge: The Earth-Moon-Sun Triangle}

\textbf{46.} During a first-quarter Moon: - Earth-Sun distance: 150
million km - Earth-Moon distance: 384,400 km - The Moon appears 90° from
the Sun (as seen from Earth)

\begin{itemize}
\item
  \begin{enumerate}
  \def\labelenumi{(\alph{enumi})}
  \tightlist
  \item
    What is the Sun-Moon distance at this moment?
  \end{enumerate}
\item
  \begin{enumerate}
  \def\labelenumi{(\alph{enumi})}
  \setcounter{enumi}{1}
  \tightlist
  \item
    What is the angle at the Moon in this triangle?
  \end{enumerate}
\item
  \begin{enumerate}
  \def\labelenumi{(\alph{enumi})}
  \setcounter{enumi}{2}
  \tightlist
  \item
    What is the angle at the Sun?
  \end{enumerate}
\end{itemize}

\begin{center}\rule{0.5\linewidth}{0.5pt}\end{center}

\end{document}
